\documentclass[12pt]{scrartcl}
\usepackage[nopatch=footnote]{microtype}

\usepackage{evan}
\usepackage{graphicx,amsmath,amssymb,setspace}

\newcommand{\gA}{\mathfrak{A}}
\newcommand{\gB}{\mathfrak{B}}
\newcommand{\cL}{\mathcal{L}}

\begin{document}

\title{MATMEK4270 Assignment 1}
\author{Sander Oppen}
\date{\today}
\maketitle

\subsection*{1.2.3. Exact solution}
We begin by finding the partial derivatives for $u$, starting with $t$.
\begin{align*}
    \frac{\partial u}{\partial t}&= e^{\imath (k_xx+k_yy)}(-\imath \omega)e^{-\imath\omega t}=-\imath\omega u\\
    \frac{\partial^2 u}{\partial t^2} &= (-\imath \omega)^2u = -\omega^2 u. 
\end{align*}
Next, for the derivatives for $x$ and $y$ (I have skipped some steps when calculating, they are the same as the previous scenario).
\begin{align*}
    \frac{\partial u}{\partial x}&= \imath k_x u\\
    \frac{\partial^2 u}{\partial x^2} &= (\imath k_x)^2u = -k_x^2 u. \\
    \frac{\partial u}{\partial y}&= \imath k_y u\\
    \frac{\partial^2 u}{\partial y^2} &= (\imath k_y)^2u = -k_y^2 u. 
\end{align*}
Putting all these into the wave equation, we get
\begin{equation}
    -\omega^2 u = c^2\big( -k_x^2u -k_y^2u\big).
\end{equation}
Since $u$ has no zeroes, we know that the wave equation holds if and only is we have $\omega = \pm c\sqrt{k_x^2+k_y^2}$.

\newpage
\subsection*{1.2.4. Dispersion coefficient}
We begin by finding the required values for $u_{ij}^n$, starting with $n+1$ and $n-1$.
\begin{align*}
    u_{ij}^{n+1}&= e^{\imath(kh(i+j)-\tilde{\omega}(n+1)\Delta t)} = e^{-\imath\tilde{\omega}\Delta t}u_{ij}^n\\
    u_{ij}^{n-1}&= e^{\imath\tilde{\omega}\Delta t}u_{ij}^n.
\end{align*}
Next, for $i\pm1$ and $j\pm 1$ (I have again skipped some steps when calculating, they are the same as the previous scenario):
\begin{align*}
    u_{i+1,j}^n&= e^{\imath kh}u_{ij}^n\\
    u_{i-1,j}^n&= e^{-\imath kh}u_{ij}^n\\
    u_{i,j+1}^n&= e^{\imath kh}u_{ij}^n\\
    u_{i,j-1}^n&= e^{-\imath kh}u_{ij}^n.
\end{align*}
Putting all these into equation (1.3) from the problem statement, we get
\begin{equation}
    \frac{e^{-\imath\tilde{\omega}\Delta t}-2+e^{\imath\tilde{\omega}\Delta t}}{\Delta t^2} = \frac{2c^2}{h^2}\bigl(e^{\imath kh}-2+e^{-\imath kh}\bigr).
\end{equation}
Using $e^{\imath a}+e^{-\imath a}=2\cos(a)$, we get
\begin{equation*}
    \frac{2\cos(\tilde{\omega}\Delta t)-2}{\Delta t^2} = \frac{2c^2}{h^2}\bigl( 2\cos(kh)-2 \bigr),
\end{equation*}
and by further using $\cos(a)-1=-2\sin^2\bigl(\frac{a}{2}\bigr)$, we get
\begin{align*}
    -\frac{4}{\Delta t^2}\sin^2\biggl(\frac{\tilde{\omega}\Delta t}{2}\biggr)&= -\frac{8c^2}{h^2}\sin^2\biggl(\frac{kh}{2}\biggr)\\
    \sin^2\biggl(\frac{\tilde{\omega}\Delta t}{2}\biggr)&= \frac{8c^2\Delta t^2}{4h^2}\sin^2\biggl(\frac{kh}{2}\biggr)\\
    \sin^2\biggl(\frac{\tilde{\omega}\Delta t}{2}\biggr)&=2C^2\sin^2(kh/2).
\end{align*}
In the last step we have used the substitution $C=\frac{c\Delta t}{h}$. For $C=\frac{1}{\sqrt{2}}$, this becomes
\begin{equation}
    \sin^2\biggl(\frac{\tilde{\omega}\Delta t}{2}\biggr) = \sin^2\biggl(\frac{kh}{2}\biggr).
\end{equation}
From this, we know that
\begin{align*}
    \frac{\tilde{\omega}\Delta t}{2}&=\frac{kh}{2}\\
    \tilde{\omega}&=\frac{kh}{\Delta t}.
\end{align*}
Going back to our constant $C$, we can express $\Delta t$ as 
\begin{equation*}
    \Delta t = \frac{h}{\sqrt{2}c}
\end{equation*}
Putting this into our previous expression for $\tilde{\omega}$, we get
\begin{equation*}
    \tilde{\omega}=\sqrt{2}ck.
\end{equation*}
Looking back at what we found in the previous problem, by setting $k_x = k_y = k$, we get that 
\begin{equation*}
    \omega=c\sqrt{2k^2}=\sqrt{2}ck = \tilde{\omega}.
\end{equation*}
\end{document}
